% Сборка: latexmk -xelatex -interaction=nonstopmode -halt-on-error ne5532_at3600l.tex
% Расчёт: .venv/bin/python scripts/calculate_unbalanced_riaa.py
\documentclass[a4paper,landscape,11pt]{article}
% TikZ загружает xcolor раньше соответствующей строки в preamble.tex.
\PassOptionsToPackage{dvipsnames}{xcolor}
\usepackage{iftex}

\ifTUTeX
\usepackage{xltxtra}

\usepackage{fontspec}
\usepackage{polyglossia}
\setmainlanguage[babelshorthands=true]{russian}
\setotherlanguage{english}

\setmainfont{Times New Roman}
\setromanfont{Times New Roman} 
\setsansfont{Arial} 
\setmonofont{Courier New} 

\newfontfamily{\cyrillicfont}{Times New Roman} 
\newfontfamily{\cyrillicfontrm}{Times New Roman}
\newfontfamily{\cyrillicfonttt}{Courier New}
\newfontfamily{\cyrillicfontsf}{Arial}

% \usepackage{unicode-math}
\usepackage{amssymb}
\usepackage{amsfonts}

\else
% For pdfTeX we must use old
% 8-bit font technologies
\usepackage[T2A]{fontenc}
\usepackage[english,russian]{babel}
\usepackage[utf8]{inputenc}
\usepackage[english, russian]{babel}

%\usepackage{newtxmath}
\usepackage{amsmath}
\usepackage{amsthm}
\usepackage{amssymb}
\usepackage{amsfonts}
%\usepackage{mathtext}
\fi

%Hyphenation rules
\usepackage{hyphenat}
\hyphenation{ма-те-ма-ти-ка вос-ста-нав-ли-вать ото-бра-жа-ют-ся сге-не-ри-ро-ван-ном сим-во-лы}
% \usepackage[english, russian]{babel}

\usepackage{graphicx}
\graphicspath{ {docs/images/} {../images/} {./images/} }

\usepackage{tikz}
\usepackage{pgfplots} % для построения графиков
\usepackage{tkz-euclide}
\pgfplotsset{compat=1.18} % версия для совместимости
\usepackage{empheq}
\usetikzlibrary{positioning, calc, intersections}

\usepackage{float}

% \usepackage{comment}

\usepackage[hidelinks]{hyperref}

\newtheorem{theorem}{Теорема}[section]
\newtheorem{corollary}{Corollary}[theorem]
\newtheorem{lemma}[theorem]{Lemma}

\addto\captionsrussian{%
  \renewcommand{\figurename}{Рис.}%
  \renewcommand{\tablename}{Табл.}%
}

\usepackage{verbatim}

\newcommand{\cg}{\textnormal{\textsl{g}}}

\usepackage[dvipsnames]{xcolor}

\usepackage[siunitx, RPvoltages, american]{circuitikzgit}
\ctikzset{
    amplifiers/fill=cyan,
    sources/fill=green!20,
    diodes/fill=white,
    resistors/fill=violet,
    grounds/scale=1.2,
}

\usepackage[margin=12mm]{geometry}
\pagestyle{empty}
\setlength{\parindent}{0pt}
\ctikzset{amplifiers/fill=cyan!8, resistors/fill=white,
  grounds/scale=0.7, bipoles/length=1.05cm}
\newcommand{\kohm}{\,\text{кОм}}
\newcommand{\ohm}{\,\text{Ом}}
\newcommand{\nf}{\,\text{нФ}}
\newcommand{\uf}{\,\text{мкФ}}

% #1: канал; #2: префикс; #3/#4/#5: вход + / вход - / выход;
% #6/#7: цвет сигнального и обратного проводов.
\newcommand{\channel}[7]{%
\begin{circuitikz}[y=0.8cm, font=\small, line width=0.6pt]
  \node[op amp, noinv input up] (U) at (10.2,0) {U1#2};
  \node[font=\footnotesize] at (10.2,-1.25) {NE5532};
  \coordinate (IN) at (0,0 |- U.+);
  \coordinate (BIAS) at (4,0 |- U.+);
  \coordinate (FB) at (7,0 |- U.-);
  \coordinate (O) at (12,0);
  \node[anchor=west,font=\bfseries] at (0,2.1) {#1 канал};
  \draw (IN) node[ocirc] {} node[above left] {#6: $#1+$}
    to[C,l=$C_{#21}\;2.2\uf$] (3,0 |- U.+)
    -- (BIAS) -- (U.+);
  \draw (BIAS) to[R,l_=$R_{#21}\;47\kohm$,*-]
    (4,-1.8) node[ground] {};
  \draw (0,-1.8) node[ocirc] {} node[above left] {#7: $#1-$}
    -- (1.7,-1.8) node[ground] {};
  \node[anchor=west,font=\footnotesize,text=black!65] at (0,-2.75)
    {Провода от AT-3600L / тонарма};
  \draw (U.-) -- (FB) -- (7,-1.3)
    to[R,l_=$R_{#24}\;91\ohm$] (7,-2.65)
    to[C,l_=$C_{#23}\;1000\uf\;\mathrm{BP}$]
    (7,-3.8) node[ground] {};
  % Низко- и среднечастотная часть RIAA: Rf || (Cf + Rz).
  \draw (FB) to[short,*-] (7,1.4)
    to[R,l=$R_{#22}\;90.9\kohm$] (12,1.4) -- (O);
  \draw (7,1.4) to[short,*-] (7,2.65)
    to[C,l=$C_{#22}\;31.5\nf$] (9.5,2.65)
    to[R,l=$R_{#23}\;10\kohm$] (12,2.65)
    to[short,-*] (12,1.4);
  \draw (U.out) -- (O)
    to[R,l=$R_{#25}\;750\ohm$,*-] (15,0)
    to[C,l=$C_{#25}\;4.7\uf$] (18,0)
    -- (21,0) node[ocirc] {} node[above] {OUT #1};
  \draw (15,0) to[C,l_=$C_{#24}\;100\nf$,*-]
    (15,-1.8) node[ground] {};
  \draw (19,0) to[R,l_=$R_{#26}\;100\kohm$,*-]
    (19,-1.8) node[ground] {};
  \draw (21,-1.8) node[ocirc] {} -- (20.5,-1.8) node[ground] {};
  \node[anchor=west,font=\footnotesize,text=black!65] at (17,-2.75)
    {К линейному входу AUX / LINE};
  \node[above left,font=\footnotesize] at (U.+) {#3};
  \node[below left,font=\footnotesize] at (U.-) {#4};
  \node[above,font=\footnotesize] at ($(U.out)+(0.2,0)$) {#5};
\end{circuitikz}%
}

\begin{document}
{\LARGE\bfseries AT-3600L → стереофонокорректор на NE5532}\par
\vspace{2mm}
Одна микросхема, два канала. Нагрузка головки $47\kohm$;
усиление около $40$ дБ на $1$ кГц; питание $\pm12$ В относительно GND.
\par\vspace{1mm}
\begin{center}
\resizebox{0.94\linewidth}{!}{\channel{L}{A}{3}{2}{1}{Белый}{Синий}}\par
\vspace{2mm}
\resizebox{0.94\linewidth}{!}{\channel{R}{B}{5}{6}{7}{Красный}{Зелёный}}
\end{center}
\vfill
\small
Все символы земли обозначают общий GND. Вывод 8 U1: $+12$ В; вывод 4: $-12$ В.
Обозначения $R_{A1}\ldots R_{A6}$ и $R_{B1}\ldots R_{B6}$ относятся к разным деталям.
Питание головке не требуется. Питание, номиналы и монтаж — на листе 2.

\newpage
{\LARGE\bfseries Питание и подключение}\par
\vspace{3mm}
\begin{minipage}[t]{0.60\linewidth}
\textbf{Двухполярный стабилизированный источник $+12$ В / GND / $-12$ В}
\par\vspace{2mm}
\resizebox{\linewidth}{!}{%
\begin{circuitikz}[font=\small, line width=0.6pt]
  \node[draw,minimum width=2cm,minimum height=2.8cm,align=center] (IC) at (12,0)
    {U1\\NE5532};
  \draw (0,2) node[ocirc] {} node[left] {$+12$ В}
    -- (12,2) -- (IC.north) node[pos=1,above right] {8};
  \draw (0,0) node[ocirc] {} node[left] {GND}
    -- (9.5,0) node[ground] {};
  \draw (0,-2) node[ocirc] {} node[left] {$-12$ В}
    -- (12,-2) -- (IC.south) node[pos=1,below right] {4};
  \draw (4,2) to[C,l=$C_{21}\;47\uf$,*-]
    (4,0) node[circ] {};
  \node at (3.6,1.5) {$+$};
  \draw (8,2) to[C,l=$C_{23}\;100\nf$,*-]
    (8,0) node[circ] {};
  \draw (4,0)
    to[C,l=$C_{22}\;47\uf$] (4,-2) node[circ] {};
  \node at (3.6,-0.65) {$+$};
  \draw (8,0)
    to[C,l=$C_{24}\;100\nf$] (8,-2) node[circ] {};
  \node[align=center,font=\footnotesize] at (8,-3.0)
    {$C_{23}$, $C_{24}$ — у выводов 8 и 4;\\питание общее для U1A и U1B};
\end{circuitikz}}
\end{minipage}\hfill
\begin{minipage}[t]{0.37\linewidth}
\textbf{Соединения}\par\vspace{1mm}
\begin{tabular}{ll}
Белый, L+ & вход левого канала\\
Синий, L− & GND левого входа\\
Красный, R+ & вход правого канала\\
Зелёный, R− & GND правого входа\\
OUT L, OUT R & центры выходных RCA\\
Земли выходов & наружные контакты RCA
\end{tabular}
\par\vspace{3mm}
Отдельный провод земли тонарма/шасси проигрывателя подвести к общей точке GND
около входных разъёмов. Возвраты L− и R− вести отдельно до этой точки.
\par\vspace{2mm}
Источник должен иметь три вывода и обеспечивать не менее 30 мА.
GND — его средняя точка; это не отрицательный вывод обычного адаптера 12 В.
\end{minipage}

\vspace{4mm}
\begin{minipage}[t]{0.48\linewidth}
\textbf{Детали и монтаж}\par\vspace{1mm}
Резисторы — металлоплёночные, 1\%, 0.125–0.25 Вт.
Для обоих каналов одинаковые номиналы.
$C_{A2}$ и $C_{B2}$ по $31.5\nf$:
каждый собрать из $30\nf+1.5\nf$, включённых параллельно;
желательно плёночные, 1\% или подобранные измерением.
$C_{A4}$, $C_{B4}$ — плёночные $100\nf$, 1\%.
\par\vspace{2mm}
$C_{A1}$, $C_{B1}$ ($2.2\uf$) и $C_{A5}$, $C_{B5}$ ($4.7\uf$) —
плёночные, от 25 В. $C_{A3}$, $C_{B3}$ — \textbf{неполярные}
электролиты (BP), $1000\uf$, от 16 В.
$C_{21}$, $C_{22}$ — электролиты 47 мкФ, 25 В; знаки «+» показаны на схеме.
$C_{23}$, $C_{24}$ — керамические 100 нФ, от 25 В.
\par\vspace{2mm}
Входные провода — короткие экранированные. Цепи обратной связи размещать у ОУ,
вдали от выходных проводов и блока питания. Металлический корпус соединить с GND
в одной точке. Входные разъёмы изолировать от корпуса.
Ёмкость кабеля входит в нагрузку головки; дополнительную входную ёмкость
подбирать с учётом кабеля и паспорта конкретной головки.
\end{minipage}\hfill
\begin{minipage}[t]{0.48\linewidth}
\textbf{Как работает коррекция}\par\vspace{1mm}
В каждом канале $R_2\parallel(R_3+1/sC_2)$ в обратной связи
задаёт постоянные времени около $3180$ и $318$ мкс.
Выходная цепь $R_5C_4=75$ мкс формирует высокочастотный спад RIAA.
Обратная связь берётся \textbf{до} $R_5$.
$C_3$ уменьшает усиление постоянного смещения; $C_5$ отделяет постоянное
напряжение на выходе.
\par\vspace{2mm}
При линейной нагрузке $47\kohm$ расчётное усиление на 1 кГц — $99.1$
($39.92$ дБ); при сигнале головки 4 мВ получается около 0.40 В.
Отклонение от идеальной RIAA, нормированной на 1 кГц, примерно от
$-0.11$ до $+0.17$ дБ в полосе 20 Гц–20 кГц.
Это расчёт с идеальным ОУ и номинальными деталями, без импеданса головки;
реальный макет требует измерения. Вход следующего устройства — $\geq47\kohm$.
\par\vspace{2mm}
Подключать выходы к \textbf{AUX/LINE}, а исходный сигнал головки подавать
на вход этой схемы. Если в проигрывателе уже есть фонокорректор, использовать
его обход / выход PHONO. При первом включении проверить шины питания и отсутствие
постоянного напряжения на выходных разъёмах.
\end{minipage}

\vfill
\footnotesize
Справочные документы:
\href{https://www.ti.com/lit/ds/symlink/ne5532.pdf}{TI: NE5532, распиновка и питание};
\href{https://www.onsemi.com/download/application-notes/pdf/and8177-d.pdf}{onsemi AND8177/D: фонокоррекция на NE5532};
\href{https://www.audio-technica.co.jp/pdf/support/AT-VM95C_UM.pdf}{Audio-Technica: цветовая маркировка выводов}.
Номиналы разделённой RIAA-цепи на листе 1 рассчитаны для данной схемы.

\newpage
\normalsize
{\LARGE\bfseries Расчёт АЧХ и проверка RIAA}\par
\vspace{3mm}
\begin{minipage}[t]{0.48\linewidth}
\raggedright
\small
\textbf{Скрипт и исходные параметры}\par
\texttt{scripts/calculate\_unbalanced\_riaa.py} рассчитывает схему листа 1
с её существующими номиналами. Общие узловые уравнения, проверка и графики —
в \texttt{scripts/riaa\_common.py}.
Номиналы находятся в секции \texttt{[unbalanced]} файла \texttt{config.toml};
\texttt{[analysis]} задаёт полосу, входной уровень и внешнюю нагрузку,
\texttt{[opamp]} — учебную модель ОУ, \texttt{[riaa]} — эталонные постоянные времени.
\par\vspace{2mm}
Зависимости NumPy и Matplotlib объявлены в \texttt{pyproject.toml}.
Установка: \texttt{.venv/bin/python -m pip install -e .}\par
Запуск из каталога проекта:\par
{\footnotesize\ttfamily .venv/bin/python scripts/calculate\_unbalanced\_riaa.py}\par
Доступны параметры \texttt{--config PATH} и \texttt{--output-dir DIR}.
Изменение TOML меняет расчёт, но автоматически не переписывает TeX-схему.
\par\vspace{2mm}
\textbf{Что считается}\par
Для $s=j2\pi f$ входной разделительный конденсатор даёт
$H_{\mathrm{in}}=sR_{A1}C_{A1}/(1+sR_{A1}C_{A1})$.
Обратная связь и нижняя ветвь:
\[
 Z_f=\left(\frac1{90.9\kohm}
 +\frac1{10\kohm+\frac1{s\,31.5\nf}}\right)^{-1},\qquad
 Z_g=91\ohm+\frac1{s\,1000\uf}.
\]
Усиление ОУ $G=1+Z_f/Z_g$.
Выходной фильтр 750 Ом / 100 нФ нагружается последовательно включённым
разделительным C и параллелью 100 кОм с входом AUX/LINE.
Скрипт учитывает весь этот делитель, включая нагрузку 47 кОм;
поэтому результат немного отличается от идеального полюса ровно 75 мкс.
\par\vspace{2mm}
Передаточная функция сравнивается с независимым узловым расчётом по Кирхгофу.
Дополнительно считается однополюсный ОУ с $A_0=10^5$, GBW = 12 МГц.
Эта модель не учитывает шум, импеданс катушки, ёмкость кабеля, допуски,
перегрузку и высокочастотную устойчивость реального NE5532.
\end{minipage}\hfill
\begin{minipage}[t]{0.49\linewidth}
\vspace{0pt}
\IfFileExists{build/balanced_riaa/unbalanced.pdf}{%
  \includegraphics[width=\linewidth]{build/balanced_riaa/unbalanced.pdf}%
}{\small График появится после запуска скрипта с исходным \texttt{config.toml}.}
\par\vspace{2mm}\small
\textbf{Результат для исходных номиналов}\par
На 1 кГц: 99.13 В/В (39.92 дБ), то есть 4 мВ дают около 0.397 В.
Отклонение от RIAA с нормировкой на 1 кГц:
от −0.107 до +0.163 дБ в полосе 20 Гц–20 кГц (идеальный ОУ).
Каждая кривая ошибки нормирована по собственному усилению на 1 кГц.
\par\vspace{2mm}
В \texttt{build/balanced\_riaa/} сохраняются:
\texttt{unbalanced.csv} — таблица по частотам;
\texttt{unbalanced.json} — итоговые числа и снимок конфигурации;
\texttt{unbalanced.png} и \texttt{unbalanced.pdf} — графики.
При другом \texttt{--output-dir} путь к рисунку в этом TeX-файле нужно изменить вручную.
\end{minipage}
\end{document}
